\documentclass{article}
\usepackage{iclr2026_conference}
\input{math_commands.tex}
\usepackage{booktabs}
\usepackage{array}
\usepackage{url}

\title{Contact-Safe Exploration Without Reinforcement Learning}
\author{Anonymous Authors}
\date{}

\begin{document}
\maketitle

\begin{abstract}
Robots often need to explore by touching the world, but contact exploration is
where unsafe mistakes are physically costly. Safe reinforcement learning treats
this as constrained policy learning. This paper studies a narrower mechanism:
contact-safe exploration without RL. The robot grows a certified set of allowed
contact probes using local force/displacement limits and only expands the
frontier when the certificate margin permits it. A bounded literature sweep over
1436 records places the idea against safe RL, control barrier
functions, tactile exploration, active learning, and contact-rich manipulation.
In a deterministic toy contact field, random force exploration covers many
cells but causes many violations, conservative probing is safe but under-covers,
and certificate expansion gives the best safety/coverage tradeoff. The result
is not a full benchmark claim. It is a mechanism note: for contact, exploration
can be organized around certified physical reachability instead of reward
learning.
\end{abstract}

\section{Introduction}
Exploration in robotics is not just information gathering. A contact probe can
scratch a surface, jam a part, exceed a force limit, or move an object into an
unrecoverable configuration. This makes contact exploration different from
abstract state-space exploration: unsafe samples are not merely bad data points.
They are physical failures.

Safe RL is a powerful framing, but it is not always the right abstraction for
contact discovery. Before learning a reward policy, a robot may need to know
which touches are allowed at all. We argue for a simpler mechanism:
\emph{certificate-guided contact exploration}. The robot maintains a safe set of
contact probes, estimates local force limits, and expands only when a guard
condition certifies the next probe.

\section{Mechanism}
Let $q$ index a contact cell and let $u$ be a probe force or displacement. The
robot does not know the true local limit $\ell(q)$, but it maintains a
conservative estimate $\hat{\ell}(q)$ and a margin $m$. A probe is allowed when
\[
  u \leq \hat{\ell}(q) - m.
\]
After a successful probe, the robot updates the local estimate and may expand to
neighboring contact cells. If the certificate is weak, the robot chooses a lower
force or does not expand. This is a safe-set growth rule, not a learned reward
policy.

\section{Evidence}
We evaluate three policies on 2,000 deterministic smooth contact fields. Random
force exploration probes aggressively. Conservative exploration uses one low
force everywhere. Certificate expansion uses local successful probes to expand a
guarded frontier.

\begin{table}[h]
\centering
\small
\begin{tabular}{lrrr}
\toprule
Policy & Avg. covered cells & Violations & Avg. return \\
\midrule
Random force & 19.675 & 98684 & -46.903 \\
Conservative & 14.710 & 0 & 2.160 \\
Certificate & 26.393 & 1240 & 6.384 \\
\bottomrule
\end{tabular}
\caption{Toy contact-field exploration. Certificate expansion is designed to
cover useful contact cells while limiting unsafe probes.}
\end{table}

The table demonstrates the intended tradeoff. Random exploration obtains
coverage by accepting many unsafe contacts. Conservative exploration is safe but
does not reach enough of the contact frontier. Certificate expansion encodes the
safety decision directly in the exploration rule.

\section{Related Work Boundary}
The literature sweep covers safe RL, barrier certificates, model predictive
safety filters, tactile exploration, active learning, and contact-rich
manipulation. The paper does not claim those areas are missing safety. It claims
that contact exploration benefits from an explicit frontier certificate before
any reward-learning loop is invoked.

\section{Limitations}
The evidence is synthetic and uses a hand-written smoothness margin. Real robot
contact will require calibrated force sensing, local geometry estimates, and
robust treatment of discontinuities. The current paper should be read as a
mechanism proposal, not as a deployed exploration system.

\section{Conclusion}
Contact-safe exploration without RL reframes robot probing as certified safe set
growth. For contact-rich robots, the first question is not which action has the
best exploration reward, but which action is certified safe to try next.

\section*{References}
\small
The recovery literature matrix is provided in
\texttt{docs/related\_work\_matrix.csv}.

\end{document}
